\documentclass[a4paper]{arrowhead}

\usepackage[yyyymmdd]{datetime}
\usepackage{etoolbox}
\usepackage[T1]{fontenc}
\usepackage[utf8]{inputenc}
\usepackage{multirow}
\usepackage{longtable}
\usepackage{amssymb}
\usepackage{float}
\usepackage{tikz}
\usetikzlibrary{positioning,arrows.meta,fit,calc,shapes.misc,shapes.geometric}

\renewcommand{\dateseparator}{-}
\setlength{\emergencystretch}{3em}
\lstset{basicstyle=\ttfamily\footnotesize,breaklines=true,columns=fullflexible,frame=single,rulecolor=\color{gray!50},backgroundcolor=\color{gray!6}}
\tikzset{
  ahsys/.style={draw=ArrowheadBlue,very thick,fill=ArrowheadBlue!6,rounded corners=2pt,minimum width=3.2cm,minimum height=1.1cm,align=center,font=\small\bfseries},
  ahext/.style={draw=ArrowheadGrey,thick,dashed,fill=gray!5,rounded corners=2pt,minimum width=3.2cm,minimum height=1.0cm,align=center,font=\small},
  ahsrv/.style={draw=ArrowheadPurple,thick,fill=ArrowheadPurple!8,rounded corners=6pt,minimum height=0.55cm,inner xsep=4pt,align=center,font=\scriptsize\ttfamily},
  ahcsm/.style={draw=ArrowheadTeal,thick,fill=ArrowheadTeal!8,rounded corners=6pt,minimum height=0.55cm,inner xsep=4pt,align=center,font=\scriptsize\ttfamily},
  ahdb/.style={draw=ArrowheadOrange,thick,fill=ArrowheadOrange!8,cylinder,shape border rotate=90,aspect=0.25,minimum height=0.9cm,align=center,font=\scriptsize},
  ahflow/.style={-{Stealth[length=2.2mm]},thick,ArrowheadGrey},
  ahlabel/.style={font=\scriptsize,fill=white,inner sep=1pt,align=center},
}

%% Special references
\newcommand{\fref}[1]{{\textcolor{ArrowheadBlue}{\hyperref[sec:functions:#1]{#1}}}}
\newcommand{\mref}[1]{{\textcolor{ArrowheadPurple}{\hyperref[sec:model:#1]{#1}}}}
\newcommand{\sref}[1]{{\textcolor{ArrowheadPurple}{#1}}}
\newrobustcmd\msubsection[2]{
  \addtocounter{subsection}{1}
  \addcontentsline{toc}{subsection}{\protect\numberline{\thesubsection}#1 \textcolor{ArrowheadPurple}{#2}}
  \subsection*{\thesubsection\quad#1 \textcolor{ArrowheadPurple}{#2}}
  \label{sec:model:#2}
}

\begin{document}

%% Arrowhead Document Properties
\ArrowheadTitle{Arrowhead Workflow SoS}

\ArrowheadType{System of Systems Description (black box)}
\ArrowheadTypeShort{SoSD}
\ArrowheadVersion{5.4.0}
\ArrowheadDate{2026-11-24}
\ArrowheadAuthor{Documentation maintenance}
\ArrowheadStatus{REVIEW DRAFT}
\ArrowheadContact{Repository maintainers}
\ArrowheadFooter{\href{www.arrowhead.eu}{www.arrowhead.eu}}
\ArrowheadSetup
%%

%% Front Page
\begin{center}
  \vspace*{1cm}
  \huge{\arrowtitle}

  \vspace*{0.2cm}
  \LARGE{\arrowtype}
  \vspace*{1cm}

  
  \vspace*{\fill}

  

  \vspace*{1cm}
  \vspace*{\fill}

  % Front Page Abstract
  \begin{abstract}
    This document describes the Arrowhead Workflow System of Systems as a black box: the workflow capability it adds to an Eclipse Arrowhead local cloud, how it is meant to be used, its properties and delimitations, and the services at its boundary. The white-box design is given in the SoSDD. Documentation edition 5.4.0.
  \end{abstract}

  \vspace*{1cm}
\end{center}

\newpage
%%

%% Table of Contents
\tableofcontents
\newpage
%%

\section{Overview}
\label{sec:overview}
Eclipse Arrowhead binds services at run time, but it cannot state the \emph{order and conditions} of service exchanges.
Putting that logic in each consumer breaks loose coupling; leaving it to a central MES brings back a bottleneck and a
single point of failure~\cite{represa2022,derhamy2018}. The Workflow SoS closes this gap. Seen from outside, it lets an
engineer \emph{design} a workflow over the microservices of a local cloud, lets products and controllers \emph{run} it
for a subject, and makes every run \emph{traceable} against a formal execution model~\cite{varga2026}, while service
binding, authorization and QoS stay with the Arrowhead core systems.

\begin{figure}[ht!]
  \centering
  \resizebox{\textwidth}{!}{
\begin{tikzpicture}
\node[draw=ArrowheadBlue,very thick,fill=ArrowheadBlue!5,rounded corners=4pt,minimum width=6.4cm,minimum height=6.0cm] (sos) at (0,0) {};
\node[font=\bfseries,ArrowheadBlue,align=center] at (0,2.4) {Workflow SoS\\[-1pt]{\scriptsize\normalfont $\ll$black box$\gg$}};
\node[font=\scriptsize,align=center,text width=5.6cm] at (0,1.2) {design, store, validate and run workflows of service exchanges in one local cloud};
\node[ahsrv,anchor=east] (s1) at (3.2,0.0) {workflowConfiguration};
\node[ahsrv,anchor=east] (s2) at (3.2,-0.8) {workflowRecipe};
\node[ahsrv,anchor=east] (s3) at (3.2,-1.6) {workflowInstance};
\node[ahsrv,anchor=east] (s4) at (3.2,-2.4) {workflowManagement};
\node[ahext] (eng) at (7.4,0.0) {Engineer / tool\\[-1pt]{\scriptsize browser console}};
\node[ahext] (prod) at (7.4,-1.6) {Product, line controller,\\ upper WorkflowManager};
\node[ahext] (op) at (7.4,-3.1) {Sysop};
\draw[ahflow] (eng.west) -- (s1.east);
\draw[ahflow] (eng.south west) -- (s2.east);
\draw[ahflow] (prod.west) -- (s3.east);
\draw[ahflow] (op.west) -- (s4.east);
\node[ahext,minimum width=4.2cm] (core) at (-8.1,2.2) {Core systems\\[-1pt]{\scriptsize ServiceRegistry, Authentication,}\\[-3pt]{\scriptsize ConsumerAuthorization, DSO}};
\node[ahext,minimum width=4.2cm] (cfg) at (-8.1,0.6) {LocalCloudConfigurator,\\ SoLocalCloudConfigurator};
\node[ahext,minimum width=4.2cm] (eq) at (-8.1,-1.0) {Equipment / application\\ systems (step services)};
\node[ahext,minimum width=4.2cm] (pr) at (-8.1,-2.6) {Intelligent Product\\[-1pt]{\scriptsize workflowProduct}};
\draw[ahflow] (sos.west |- core) -- (core.east) node[midway,above,font=\scriptsize,align=center] {discovery, identity,\\ authorization, orchestration};
\draw[ahflow] (sos.west |- cfg) -- (cfg.east) node[midway,above,font=\scriptsize] {designs};
\draw[ahflow] (sos.west |- eq) -- (eq.east) node[midway,above,font=\scriptsize] {invoke};
\draw[ahflow] (sos.west |- pr) -- (pr.east) node[midway,above,font=\scriptsize] {recipe, results};
\node[ahext] (ch) at (0,-4.6) {Choreographer\\[-1pt]{\scriptsize (compatibility)}};
\draw[ahflow,<->] (sos.south) -- (ch.north) node[midway,right,font=\scriptsize,align=left] {executor interface,\\ compiled plans};
\end{tikzpicture}
  }
  \caption{The Workflow SoS as a black box: services at its boundary (right) and the systems it relies on (left and below).}
  \label{fig:blackbox}
\end{figure}

The rest of this document is organized as follows. Section~\ref{sec:prior_art} references prior art,
Section~\ref{sec:use} the intended use, Section~\ref{sec:properties} the properties, Section~\ref{sec:delimitations}
the delimitations, Section~\ref{sec:services} the services at the boundary and Section~\ref{sec:security} security.

\newpage
\subsection{Significant Prior Art}
\label{sec:prior_art}
\begin{itemize}
\item Derhamy~\cite{derhamy2018,derhamy2018workflow}: the split into Intelligent Product, Workflow Manager and Workflow
Executor; the Workflow Manager as a cloneable blackboard per work centre; the product as final arbiter; the design
principles loose coupling, searchable, run-time binding, autonomy, specialisation, data at its source and first-person
perspective.
\item Garcia Represa~\cite{represa2022}: the ordering gap in microservice SoS; workflow management as a separate system;
coexistence of orchestration-style and choreography-style workflow management; workflows are selected, implemented
and executed.
\item Varga~\cite{varga2026}: the formal recipe $W$, its well-formedness (W1)--(W4), the session state and transition
rules, and the guarantees Theorem~1 (at most $3|S| + 2|A|$ transitions), Theorem~2 (progress) and Proposition~1
(exclusive tool use).
\item The Arrowhead Choreographer (4.6 to 5.4): plans, sessions and the executor contract, with which the SoS
coexists~\cite{delsing2017}.
\end{itemize}

\subsection{How This SoS Is Meant to Be Used}
\label{sec:use}
\begin{enumerate}
\item \textbf{Design.} An engineer opens the WorkflowConfigurator console, picks a configured local cloud, draws a
workflow with conditions between its microservices, validates it and stores it.
\item \textbf{Orchestrated run (push).} A line controller starts a session; the SoS selects an executor for each
enabled step, which has the provider orchestrated and invokes it.
\item \textbf{Blackboard run (pull).} Workstation executors claim the steps they can do; suitable for plug-and-produce
stations and intermittent links.
\item \textbf{Product-driven run.} A product arriving at a workstation hands over its recipe fragment and receives the
results when the session closes.
\item \textbf{Shared tools.} Several lines share a tool under expiring leases, so a failed line does not block the others.
\item \textbf{Legacy plans.} Recipes compile to Choreographer plans, and executors serve the Choreographer's executor
interface, so existing Choreographer installations keep working.
\end{enumerate}

\subsection{SoS functionalities and properties}
\label{sec:properties}
\subsubsection{Functional properties of the SoS}
\begin{itemize}
\item Workflow design over the services of the live cloud, LocalCloudConfigurator designs and SoLC member clouds.
\item Recipe validation with 30 coded checks and a verdict (ACCEPT, REVISE, CLARIFY, REJECT); immutable versioned storage.
\item Execution with the eight transition rules of the formal model, push and pull dispatch, guarded transitions,
retries, timeouts, aborts and leases; full traces.
\item Coexistence with the Choreographer in both directions.
\end{itemize}
\subsubsection{Configuration of SoS properties}
Per deployment: the workstation identifier of the WorkflowManager, the capabilities, services, workstation and mode of
each executor, the dispatch default, timeouts and lease durations, and the ConsumerAuthorization grants between the
systems and towards the equipment services. Per workflow: the recipe.
\subsubsection{Data stored by the individual microsystems}

{\small
\begin{longtable}{|p{4.0cm}|p{11.8cm}|}

\hline
\rowcolor{gray!33} System & Data \\ \hline
\endhead
WorkflowManager & Recipes (immutable), sessions (state $z$ and trace), work items, leases, worklog: database \texttt{ah\_workflow}, embedded in the WorkflowManager and reachable only through its services. \\ \hline
WorkflowExecutor & No workflow data; work items in progress in memory. \\ \hline
WorkflowConfigurator & No workflow data; console sessions in memory. \\ \hline
\end{longtable}
}

\subsubsection{Non functional properties}
\begin{itemize}
\item \emph{Security:} mutual TLS, identity tokens bound to certificates, ConsumerAuthorization on every operation,
object-level checks (Section~\ref{sec:security}).
\item \emph{Safety and predictability:} for well-formed recipes without extensions, bounded and progressing runs
(Theorems~1--2); failure is absorbing; every run is checked against the model.
\item \emph{Loose coupling:} recipes name service definitions and capabilities, never endpoints; providers are bound
per step and per session, so new equipment is used without changing the workflow (plug-and-produce).
\item \emph{Autonomy:} a workstation cloud with its core systems, one WorkflowManager, its executors and equipment can
run without a factory connection.
\end{itemize}
\subsubsection{Stateful or stateless}
Stateful in the WorkflowManager (recipes and sessions survive restarts); executors and the configurator are stateless.

\subsection{Important Delimitations}
\label{sec:delimitations}
\begin{itemize}
\item One local cloud: federation of WorkflowManagers across clouds and the disconnected working copy are designed but
not implemented.
\item No planning, scheduling or supply-chain functions; no equipment services.
\item EventHandler-based (\texttt{event}) dispatch and workflow events are not implemented.
\item Equipment adapters for HTTP/HTTPS JSON only.
\item Extensions (rebind, retry, timeout, abort, guards, compensation) are outside the proven model; sessions using them are
flagged.
\item Java only: the Go and Rust implementations do not exist, so the SoS is not yet interoperable across the three
implementations.
\end{itemize}

\newpage
\section{Services}
\label{sec:services}
Only the services below cross the boundary of the SoS. The services between its constituent systems
(\texttt{workflow\allowbreak{}Execution}, \texttt{workflow\allowbreak{}Work\allowbreak{}Item}, \texttt{workflow\allowbreak{}Context}, \texttt{workflowLease}) are internal and
described in the SoSDD, although they are ordinary Arrowhead services that other systems could use.

\subsection{Produced services}

{\small
\begin{longtable}{|p{4.6cm}|p{3.6cm}|p{7.6cm}|}

\hline
\rowcolor{gray!33} Service & Provider & Consumers and purpose \\ \hline
\endhead
\texttt{workflow\allowbreak{}Configuration} & WorkflowConfigurator & Engineers and engineering tools: clouds, services, executors; validate and store workflows. \\ \hline
\texttt{workflow\allowbreak{}Recipe} & WorkflowManager & Engineering tools, upper WorkflowManagers: recipe catalogue, compile to Choreographer plans. \\ \hline
\texttt{workflow\allowbreak{}Instance} & WorkflowManager & Products, line controllers, upper WorkflowManagers: start, observe and abort sessions. \\ \hline
\texttt{workflow\allowbreak{}Management} & WorkflowManager & Sysop: worklog, executors, work items. \\ \hline
\texttt{general\allowbreak{}Management} & all three & Sysop: logs and configuration. \\ \hline
Choreographer executor interface & WorkflowExecutor & The Choreographer (compatibility, not registered). \\ \hline
\end{longtable}
}


\subsection{Consumed services}

{\small
\begin{longtable}{|p{5.6cm}|p{4.2cm}|p{6.0cm}|}

\hline
\rowcolor{gray!33} Service & Provider & Purpose \\ \hline
\endhead
\texttt{service\allowbreak{}Discovery}, \texttt{service\allowbreak{}Registry\allowbreak{}Management} & ServiceRegistry & Registration, executor discovery, cloud listing. \\ \hline
\texttt{identity} & Authentication & Identity tokens of the systems. \\ \hline
\texttt{authorization} & ConsumerAuthorization & Provider-side checks of every operation. \\ \hline
\texttt{service\allowbreak{}Orchestration} & Dynamic\allowbreak{}Service\allowbreak{}Orchestration & Provider of each step, with the executor as requester. \\ \hline
step services & Equipment and application systems & The work of each step. \\ \hline
\texttt{workflow\allowbreak{}Product} & Intelligent Product & Recipe retrieval and result storage. \\ \hline
\texttt{local\allowbreak{}Cloud\allowbreak{}Configuration}, \texttt{so\allowbreak{}Local\allowbreak{}Cloud\allowbreak{}Configuration} & LocalCloudConfigurator, SoLocalCloudConfigurator & Designed clouds for the editor. \\ \hline
\texttt{choreographer\allowbreak{}Executor} & Choreographer & Registration and step notification (compatibility). \\ \hline
\end{longtable}
}


\newpage
\section{Security}
\label{sec:security}
\begin{itemize}
\item The SoS runs only in Arrowhead secure mode in a cloud: every system has a certificate of the local cloud CA, all
traffic uses TLS with client authentication, and every request carries an identity token bound to the caller's
certificate.
\item Non-compliant certificates and certificates on the CRL are refused.
\end{itemize}
\subsection{Security model}
\begin{itemize}
\item \textbf{Protocol:} HTTPS (TLS 1.3/1.2) with JSON; interface policy \texttt{CERT\_AUTH}.
\item \textbf{Data protection:} TLS on all links; each system's database is embedded in its own process, so no database link or port exists; secrets from container secrets.
\item \textbf{System authentication:} X.509 client certificates plus Authentication identity tokens.
\item \textbf{Authorization of produced services:} ConsumerAuthorization grant per service operation at the provider;
management services sysop only.
\item \textbf{Authorization towards equipment:} each step's provider is orchestrated with the \emph{executor} as
requester, so the equipment authorizes the executor, never the WorkflowManager or the session owner.
\item \textbf{Object level:} only the session owner may abort, only the work-item holder may report or use its context,
only the configured WorkflowManager may start work on an executor.
\item \textbf{Human access:} the WorkflowConfigurator console is a separate loopback connector with a password (scrypt
hash), session cookie, Host/Origin checks and lockout.
\end{itemize}

\newpage
\bibliographystyle{IEEEtran}
\bibliography{bibliography}

\newpage

\section{Revision History}
\subsection{Amendments}

\noindent\begin{tabularx}{\textwidth}{| p{1cm} | p{2.2cm} | p{1.6cm} | X | p{3.6cm} |} \hline
\rowcolor{gray!33} No. & Date & Version & Subject of Amendments & Author \\ \hline
1 & 2026-11-23 & \arrowversion & Initial edition: black-box description of the Workflow SoS, written from the Java reference implementation (ah5-Workflow-java). & Documentation maintenance \\ \hline
2 & 2026-11-24 & \arrowversion & Embedded database instead of a separate database server; data ownership. & Documentation maintenance \\ \hline
\end{tabularx}

\subsection{Quality Assurance}

\noindent\begin{tabularx}{\textwidth}{| p{1cm} | p{2.2cm} | p{1.6cm} | X |} \hline
\rowcolor{gray!33} No. & Date & Version & Approved by \\ \hline
1 & & \arrowversion & Pending review \\ \hline
\end{tabularx}

\end{document}
